% Part: incompleteness
% Chapter: incompleteness-provability
% Section: lob-thm

\documentclass[../../../include/open-logic-section]{subfiles}

\begin{document}

\olfileid{inc}{inp}{tar}

\olsection{Ora Bisane Ngedefinisikake Bebeneran}

Pangerten \emph{definabilitas} gumantung marang anane semantik formal
kanggo basa aritmetika. Kita wis njlentrehake sawijining himpunan
formula lan ukara ing basa aritmetika. ``Interpretasi kang
dikarepake'' yaiku maca ukara-ukara kasebut minangka pratelan
babagan wilangan asli, lan pratelan mau bisa bener utawa salah.
Upamana $\Struct{N}$ iku !!{structure} kanthi domain $\Nat$ lan
interpretasi standar kanggo simbol-simbol ing basa aritmetika.
Mula $\Sat{N}{!A}$ tegese ``$!A$ bener ing interpretasi standar.''

\begin{defn}
Relasi $R(x_1,\dots,x_k)$ antarane wilangan asli \emph{bisa
didefinisikake} ing $\Struct{N}$ yen lan mung yen ana formula
$!A(x_1,\dots,x_k)$ ing basa aritmetika saengga kanggo saben
$n_1,\dots,n_k$, $R(n_1,\dots,n_k)$ yen lan mung yen
$\Sat{N}{!A(\num n_1,\dots,\num n_k)}$.
\end{defn}

Kanthi tembung liya, sawijining relasi bisa didefinisikake ing
$\Struct{N}$ yen lan mung yen bisa direpresentasikake ing teori
$\Th{TA}$, ing ngendi $\Th{TA} = \Setabs{!A}{\Sat{N}{!A}}$
minangka himpunan ukara aritmetika kang bener. (Yen perkara iki durung
cetha, balenana maca definisi-definisie lan priksa dhewe
kenapa mangkono.)

\begin{lem}
Saben relasi komputabel bisa didefinisikake ing~$\Struct{N}$.
\end{lem}

\begin{proof}
Gampang dipriksa yen formula kang merepresentasikake sawijining
relasi ing $\Th{Q}$ netepake relasi kang padha ing $\Struct{N}$.
\end{proof}

Saiki kita bisa takon apa kosok baline uga bener: apa saben relasi
kang bisa didefinisikake ing~$\Struct{N}$ iku komputabel?
Wangsulane ora. Tuladhane:

\begin{lem}
Relasi mandheg bisa didefinisikake ing $\Struct{N}$.
\end{lem}

\begin{proof}
Elinga yen teorema wujud normal Kleene nyatakake: saben fungsi
komputabel parsial~$f$ nduweni indeks~$e$ saengga $f(x) =
U(\umin{s}{T(e,x,s)})$ kanggo saben $x \in \Nat$, lan $U$ uga
$T$ iku rekursif primitif, mula kalorone total. Dadi, $f(x)$
nduweni nilai (yaiku komputasine mandheg) yen lan mung yen ana~$s$
saengga $T(e,x,s)$ bener.

Saiki upamana $H$ iku relasi mandheg, yaiku
\[
H = \Setabs{\tuple{e,x}}{\lexists[s][T(e, x, s)]}.
\]
Upamana $!D_T$ netepake $T$ ing $\Struct{N}$. Mula
\[
H = \Setabs{\tuple{e,x}}{\Sat{N}{\lexists[s][!D_T(\num e, \num x, s)]}},
\]
Dadi $\lexists[s][!D_T(z, x, s)]$ netepake~$H$ ing
$\Struct{N}$.
\end{proof}

\begin{prob}
Tuduhna yen $Q(n) \defiff n \in \Setabs{\Gn{!A}}{\Th{Q} \Proves !A}$
  bisa didefinisikake ing aritmetika.
\end{prob}

Kepiye karo $\Th{TA}$ dhewe? Apa teori iki bisa didefinisikake ing
aritmetika? Tegese: apa himpunan
$\Setabs{\Gn{!A}}{\Sat{N}{!A}}$ bisa didefinisikake ing
aritmetika? Teorema Tarski menehi wangsulan ora.

\begin{thm}
\ollabel{thm:tarski}
Himpunan !!{sentence}s aritmetika kang bener ora bisa
didefinisikake ing aritmetika.
\end{thm}

\begin{proof} 
Upamana $!D(x)$ netepake himpunan kasebut, yaiku
$\Sat{N}{!A}$ yen lan mung yen $\Sat{N}{!D(\gn{!A})}$.
Miturut lema titik tetep, ana formula $!A$ saengga
$\Th{Q} \Proves !A \liff \lnot !D(\gn{!A})$, lan mula
$\Sat{N}{!A \liff \lnot !D(\gn{!A})}$. Nanging banjur
$\Sat{N}{!A}$ yen lan mung yen
$\Sat{N}{\lnot !D(\gn{!A})}$, kosok balen karo anggepan
yen $!D(y)$ netepake himpunan pratelan aritmetika kang bener.
\end{proof}

Tarski ngetrapake analisis iki marang pangerten filsafati ngenani
bebeneran kang luwih umum. Kanggo basa $L$ sembarang, miturut
Tarski, pangerten bebeneran kang nyukupi kanggo $L$ kudu netepi,
kanggo saben ukara $X$,
\begin{quote}
`$X$' bener yen lan mung yen $X$.
\end{quote}
Tuladha Tarski saka basa Inggris kang asring dikutip, yen
diterjemahake, yaiku ukara
\begin{quote}
`Salju iku putih' bener yen lan mung yen salju iku putih.
\end{quote}
Nanging, kanggo saben basa kang cukup kuwat kanggo
merepresentasikake fungsi diagonal, lan kanggo saben predikat
linguistik $T(x)$, kita bisa mbangun ukara $X$ kang nyukupi
``$X$ yen lan mung yen ora $T(\text{`$X$'})$.''
Amarga kita ora kepengin predikat bebeneran nyatakake manawa ana
ukara kang bener lan salah bebarengan, Tarski ndudut yen predikat
bebeneran kanggo kabeh ukara ing sawijining basa ora bisa
ditetepake tanpa metu, kanthi cara tartamtu, saka watesing basa
kasebut. Kanthi tembung liya, predikat bebeneran kanggo sawijining
basa ora bisa didefinisikake ing basa iku dhewe.

\end{document}
